\clearpage
\section{Terminology}

\subsection{Address Values}

An \emph{effective address (EA)} is the address value or operand designator produced by an addressing mode before
segment pre-translation. A memory EA produces an address, while register and immediate EAs name non-memory operands.

A \emph{pre-segment virtual address} is the ABI-visible address before segment pre-translation. Ordinary pointers,
ELF symbol values, near CALL targets, and near JMP targets use this address space unless a rule says otherwise.

A \emph{linear address} is the result of segment pre-translation. The page-table walker consumes it when paging is
enabled; otherwise the memory system uses it directly.

A \emph{physical address} is the byte address produced by page-table translation for ordinary byte-addressed memory.
With paging disabled, the linear address is used directly instead.

A \emph{memory-system address} is the value presented to the memory subsystem after every active architectural
translation stage. Depending on the final access attributes, the subsystem may interpret it as an ordinary byte
address or an externally acknowledged bus-sized access.

\emph{Address generation} is the instruction-local calculation that combines registers, displacement, index,
scale, and immediate fields into an effective address. It does not by itself read memory.

\subsection{Segmentation and Paging}

\emph{Segment pre-translation} is the stage between effective-address calculation and paging. A disabled segment
passes the EA through. A translated window checks an offset and adds its base, while a bounds-only window checks an
already-linear address without adding the base.

A \emph{segment register image} is the 64-bit architectural value containing the segment base page, size exponent,
size mantissa, and bounds-only enable bit.

A \emph{disabled segment} has a zero size mantissa. It performs no segment-window check or base addition and passes
the effective address through as the linear address.

A \emph{translated segment} is enabled with bounds-only mode clear. It checks the effective address as an offset
within the segment span and then adds the segment base.

A \emph{bounds-only segment} is enabled with bounds-only mode set. It treats the effective address as already linear
and checks only that it lies inside the segment window.

\emph{Page-table translation} is the optional \texttt{PTCR.PE}-controlled stage that maps a canonical linear
address to a physical frame and applies PTE permissions, cache policy, and addressing type.

A \emph{PFN} is a physical frame number. A leaf PTE combines its PFN with the page offset to form the physical
address.

\subsection{Control Flow and Linkage}

A \emph{near control transfer} is a CALL, RET, JMP, or conditional branch that changes only PC within the current
code-segment context.

A \emph{segment-changing control transfer} updates CS and PC in one architectural commit. Long jumps, long calls,
and long returns belong to this category.

A \emph{linkage domain} is an ABI-visible region in which ordinary code addresses name near-call entry points and
function pointers carry no explicit CS value.

A \emph{veneer} is a small runtime- or linker-generated near-call entry point that hides a more complex transfer
sequence from ordinary callers.

A \emph{far data address} is a C data-address pair containing a canonical pre-segment virtual address and a
canonical disabled or bounds-only segment image. The optional C far-data profile gives the pair a 128-bit language
ABI representation; the base ISA still permits translated segments internally and does not make the pair an
unforgeable capability.

A \emph{far-data window} is an architectural segment register temporarily loaded from a far data address so an
explicit-segment memory operation can perform the access. The optional C far-data profile assigns this role to GS4.

\subsection{Architectural State}

A \emph{general-purpose register} is one of the sixteen 64-bit Rn integer registers, R0 through R15. Ordinary
integer register fields name this unified class.

An \emph{Rn register} is a general-purpose integer register selected by a four-bit Rn field. It may hold integers,
addresses, counters, selectors, or temporary data according to the instruction form.

The \emph{stack pointer register} is SP. It lies outside the four-bit Rn namespace and uses SS as the fixed segment
for SP-relative addressing.

A \emph{segment register operand} is a 64-bit segment register exposed through the S operand class. That class names
the segment registers in the SREG selector table, not FLAGS, STATUS, PC, or SP.

A \emph{state register} is named architectural state such as FLAGS or STATUS. Dedicated read and write instructions
access these registers instead of the S operand class.

A \emph{control register} is a privileged register exposed through the CR operand class, such as PTCR, ASCR, ECR,
SPC, SCS, and SDS.

\subsection{Compatibility}

\emph{Reserved} means held for a future architectural definition, so software does not depend on the current value
or behavior. \emph{Must be zero} requires software to write zero, and \emph{read-as-zero} means reads return zero
while the field remains reserved. \emph{Ignored} means hardware accepts the field without using it for the current
operation. \emph{Preserved} requires software to retain the previous value when rewriting the containing object.

\emph{Software-defined} state is explicitly reserved for software use and is not consumed by a future extension
without redefining the containing structure. \emph{Implementation-defined} behavior is selected by the
implementation and published through CPUID, platform documentation, or firmware. \emph{Illegal} use of an opcode,
effective-address form, selector, or operation raises the exception named by the defining rule.
